% sections/spacetime.tex

\section{Spacetime from the Peirce complement}
\label{sec:spacetime}

The Peirce complement $\Vzero = \htwo{O}$ is a 10-dimensional Jordan
subalgebra. We show that the C*-bottleneck projects it to a
4-dimensional Minkowski space $\htwo{C}_u \cong \mathbb{R}^{3,1}$,
and that the Kantor-Koecher-Tits construction on $\htwo{C}_u$ gives
the 4-dimensional conformal algebra $\mathfrak{so}(4,2)$.

\subsection{The projection $\pi_u$}

The observer's complex structure $u = e_7 \in S^6 \subset \im(\mathbb{O})$
determines a complex subalgebra $\mathbb{C}_u = \mathbb{R} \oplus
\mathbb{R} u \subset \mathbb{O}$. The projection $\pi_u : \htwo{O} \to
\htwo{C}_u$ acts on an element $Y \in \htwo{O}$ by projecting the
off-diagonal octonion entry $x_1$ onto $\mathbb{C}_u$:
\begin{equation}
\label{eq:pi-u}
\pi_u \begin{pmatrix} \beta & x_1 \\ \bar{x}_1 & \gamma \end{pmatrix}
= \begin{pmatrix} \beta & \mathrm{proj}_{\mathbb{C}_u}(x_1) \\
\overline{\mathrm{proj}_{\mathbb{C}_u}(x_1)} & \gamma \end{pmatrix}.
\end{equation}

\begin{proposition}
\label{prop:pi-u}
The projection $\pi_u$ has the following properties:
\begin{enumerate}[label=(\alph*)]
\item $\pi_u$ is idempotent, with 4-dimensional image
  $\htwo{C}_u$ and 6-dimensional kernel $W$ consisting of elements of
  $\htwo{O}$ with zero diagonal entries and off-diagonal octonion
  component in $\mathrm{span}\{e_1, \ldots, e_6\}$ (the imaginary
  octonion units orthogonal to $u$). That is,
  $W = \bigl\{\bigl(\begin{smallmatrix} 0 & w \\ \bar{w} & 0
  \end{smallmatrix}\bigr) : w \in \mathrm{span}\{e_1, \ldots, e_6\}
  \bigr\}$.

\item The Gram matrix of $\det_2$ restricted to $\htwo{C}_u$ has
  eigenvalues $\{+1, -1, -1, -1\}$: Minkowski signature $(1,3)$.

\item $\Vzero$ closure under the Jordan product is a standard
  consequence of the Peirce multiplication rules; we have verified
  this computationally as a consistency check (all $55$ basis pair
  products have zero $\Vhalf$-leakage).

\item $\pi_u$ is \emph{not} a Jordan homomorphism. The failure is
  characterized by
  \begin{equation}
  \label{eq:delta}
  \Delta(A, B) \coloneqq \pi_u(\jp{A}{B}) - \jp{\pi_u(A)}{\pi_u(B)}
  = \langle w_A, w_B \rangle_W \cdot I_2,
  \end{equation}
  where $w_A = \mathrm{proj}_W(x_{1,A})$ is the $W$-component of $A$.
  In particular, $\Delta$ vanishes on $\htwo{C}_u$ (where $w = 0$) and
  is nonzero only for elements with components in the internal
  $W$-directions.
\end{enumerate}
\end{proposition}

\begin{remark}
\label{rem:pi-u-non-hom}
No step in the derivation chain requires $\pi_u$ to be a Jordan
homomorphism. The GST Lagrangian is defined on the full 27-dimensional
$\hthree$ via $d_{IJK}$; $\pi_u$ enters only to identify which 4
dimensions constitute spacetime. The failure $\Delta$ of
Eq.~(\ref{eq:delta}) is confined to cross-terms with internal
$W$-directions and does not affect the Peirce block structure of
$d_{IJK}$.
\end{remark}

The 10-dimensional $\Vzero$ therefore splits cleanly as
\begin{equation}
\label{eq:v0-split}
\Vzero = \underbrace{\htwo{C}_u}_{4\text{-dim, spacetime}}
  \oplus \underbrace{W}_{6\text{-dim, internal}},
\end{equation}
where $\htwo{C}_u$ carries Minkowski signature from $\det_2$ and $W$ is
the kernel of $\pi_u$. The explicit Minkowski coordinates on $\htwo{C}_u$
are
\begin{equation}
\label{eq:minkowski-coords}
y^0 = \tfrac{1}{2}(\beta + \gamma), \quad
y^3 = \tfrac{1}{2}(\beta - \gamma), \quad
y^1 = \re(x_1), \quad
y^2 = \langle x_1, e_7 \rangle,
\end{equation}
where $x_1 \in \mathbb{O}$ is the off-diagonal octonion entry of $Y \in
\htwo{O}$, and $\det_2 = (y^0)^2 - (y^1)^2 - (y^2)^2 - (y^3)^2$.

\subsection{Minkowski structure from $\Vhalf$ products}

The Peirce multiplication rule $\Vhalf \circ \Vhalf \subset \Vone
\oplus \Vzero$ provides additional structure. The $\Vzero$-component
of the $\Vhalf \times \Vhalf$ product, projected through $\pi_u$,
yields four $16 \times 16$ Minkowski matrices $M_\mu$ satisfying
\begin{equation}
\label{eq:clifford}
\{M_i, M_j\} = \tfrac{1}{2}\,\delta_{ij}\, I_{16},
\quad i, j \in \{1, 2, 3\},
\qquad M_0 = \tfrac{1}{2}\, I_{16}.
\end{equation}
The rescaled matrices $\gamma_i = 2M_i$ satisfy $\{\gamma_i, \gamma_j\}
= 2\delta_{ij} I_{16}$, generating $\Cl{3,0}$ on $\mathbb{R}^{16}$.
This Clifford algebra structure means that fermion bilinears in $\Vhalf$
naturally generate the spatial rotation algebra of $\htwo{C}_u$, connecting
the matter sector to the spacetime sector at the algebraic level.

\subsection{Conformal algebra}
\label{sec:conformal}

The Kantor-Koecher-Tits (KKT) construction~\cite{Tits1962,Kantor1964,Koecher1967}
builds a Lie algebra from any Jordan algebra $J$ (see~\cite{McCrimmon2004}
for a modern treatment):
\begin{equation}
\label{eq:kkt}
\KKT(J) = J^- \oplus \Str_0(J) \oplus J^+,
\end{equation}
where $J^+$ (translations) and $J^-$ (special conformal transformations)
are copies of $J$ as a vector space, and $\Str_0(J)$ is the structure
algebra (derivations plus multiplication operators).

\begin{proposition}
\label{prop:kkt}
For the observer's spacetime $\htwo{C}_u$:
\begin{enumerate}[label=(\alph*)]
\item $\KKT(\htwo{C}_u) = \mathfrak{so}(4,2)$, the 15-dimensional
  4d conformal algebra.

\item The Killing form of $\KKT(\htwo{C}_u)$ has signature $(8, 7)$.

\item The derivation algebra $\mathrm{Der}(\htwo{C}_u)$ has dimension~3
  and generates the rotation subalgebra $\mathfrak{so}(3)$.

\item The structure algebra $\Str_0(\htwo{C}_u)$ has dimension~7 and
  contains both rotations and Lorentz boosts.

\item Every rank-1 idempotent in $\hthree$ gives the same conformal
  algebra (observer independence, via $\Ffour$ transitivity).
\end{enumerate}
\end{proposition}

The conformal algebra $\mathfrak{so}(4,2)$ contains the Poincar\'e
algebra as a subalgebra: translations ($J^+$), Lorentz transformations
(within $\Str_0$), and dilatations (the grading element). The emergence
of $\mathfrak{so}(4,2)$ from the Peirce complement is purely algebraic:
no manifold, metric, or differential structure is assumed. The KKT
construction builds these symmetries from the Jordan product alone.

\subsection{Uniqueness of the spacetime}

Not every 4-dimensional subspace of $\Vzero = \htwo{O}$ gives a
spacetime with the correct conformal algebra.

\begin{proposition}
\label{prop:uniqueness}
Let $J_4 \subset \htwo{O}$ be a 4-dimensional unital Jordan subalgebra
(containing the identity).
\begin{enumerate}[label=(\alph*)]
\item All 4-dimensional unital Jordan subalgebras of $\htwo{O}$
  (containing the identity) are isomorphic to $\mathrm{JSpin}(3)$
  (a spin factor), parametrized by the Grassmannian $\mathrm{Gr}(3, 9)$.

\item $\dim \KKT(\mathrm{JSpin}(n)) = (n+3)(n+2)/2$.
  Only $n = 3$ gives $\dim = 15 = \dim\,\mathfrak{so}(4,2)$.

\item The complex structure $u \in S^6$ uniquely selects $\htwo{C}_u$
  as the $\pi_u$-image. Different choices of $u$ are related by $G_2$
  automorphisms of $\mathbb{O}$ and give isomorphic spacetimes.
\end{enumerate}
\end{proposition}

The choice of complex structure $u$ is therefore observer-independent:
no physical content depends on which $u \in S^6$ is selected, since
$G_2$ transitivity gives isomorphic Peirce decompositions for every
choice. The full orbit is $G_2/\SU{3} \cong S^6$, and $G_2$ acts as an
internal symmetry relating equivalent descriptions.

\subsection{Stabilizer structure}
\label{sec:stabilizer}

The choice of complex structure $u = e_7$ selects a $3 + 6$ split of
the 9 traceless directions of the spin factor $\htwo{O}$: 3 lie in the
traceless subspace of $\htwo{C}_u$ (the spatial directions $y^1, y^2,
y^3$) and 6 span the complement $W$ (internal).
This split determines the internal symmetries visible to the observer.

\begin{proposition}
\label{prop:stabilizer}
The subgroup of $\Spin{9}$ preserving the $3 + 6$ split induced by
$u = e_7$ has Lie algebra $\mathfrak{so}(3) \oplus \mathfrak{so}(6)$,
dimension~18:
\begin{enumerate}[label=(\alph*)]
\item $\mathfrak{so}(3)$ acts on the spacetime indices
  $\{0, 1, 2, 3\}$ of $\htwo{C}_u$, with all three generators
  satisfying
  $\eta L_{\mathrm{Mink}} + L_{\mathrm{Mink}}^T \eta = 0$.
  This is the rotation subalgebra of the Lorentz group.

\item $\mathfrak{so}(6) \cong \mathfrak{su}(4)$ acts on the internal
  $W$-directions. A rank obstruction prevents
  $\mathfrak{g}_{\mathrm{SM}}$ (rank~4) from embedding in
  $\mathfrak{su}(4)$ (rank~3); the reduction to the Standard Model
  gauge algebra requires the Todorov-Drenska intersection mechanism
  inside $\Ffour$ (see gap~G8).

\item The cross-brackets $[\mathfrak{so}(3), \mathfrak{so}(6)] = 0$:
  spacetime rotations and internal symmetries commute.

\item $\pi_u$ is equivariant under all 18 generators
  (verified to precision $5 \times 10^{-17}$).
\end{enumerate}
\end{proposition}

The compact $\Spin{9}$ cannot contain the non-compact boost generators
of $\SO{3,1}$. This is standard: gauge symmetries (compact) and
spacetime symmetries (non-compact) belong to different sectors. Boosts
arise from the spacetime metric $\det_2$ on $\htwo{C}_u$ (which has
$\SL{2,\mathbb{C}}$ isometry), not from the internal symmetry group.
The rotation subalgebra $\mathfrak{so}(3)$ is the maximal compact
subalgebra of $\mathfrak{so}(3,1)$.
